% ── Abstract ──

\chapter*{Abstract}
\addcontentsline{toc}{chapter}{Abstract}

The growing adoption of Internet of Medical Things (IoMT) devices in healthcare environments has introduced significant cybersecurity challenges. These devices---ranging from wearable sensors to smart infusion pumps---transmit sensitive patient data across networks that are increasingly targeted by adversaries. While machine learning (ML)-based intrusion detection systems (IDSs) have shown promise in detecting network intrusions, their effectiveness is fundamentally constrained by the quality and representativeness of the training data.

A critical limitation in current research is the reliance on individual, isolated datasets that capture only a narrow slice of the attack landscape. Models trained on such data often fail to generalize across different network topologies, device configurations, and attack vectors. This thesis addresses this gap by proposing and executing a systematic methodology for integrating multiple heterogeneous IoMT network traffic datasets into a single, unified, ML-ready dataset suitable for training robust IDS models.

The methodology centers on the use of MIOTTA-NPT, a specialized network preprocessing tool that extracts a standardized set of statistical features from raw packet captures using a configurable window-based aggregation approach. By processing disparate datasets through a common feature extraction pipeline, the fundamental challenge of feature heterogeneity across datasets is resolved at the source. The integration pipeline further incorporates label harmonization across three granularity levels, data quality assurance procedures, feature scaling, and dataset provenance tracking to enable rigorous cross-domain evaluation.

The resulting unified dataset preserves the diversity of the constituent sources---including varied attack types, network protocols, and device behaviors---while ensuring structural consistency required for effective model training. The Phase~1 release (8,011,534 window records with harmonized labels and scaled features) is publicly available in the CORA Research Data Repository. Three tree-based classifiers---Decision Tree, Random Forest, and XGBoost---were trained under both single-source and unified training regimes and validated on held-out test partitions as well as on HospNet26, a real hospital network traffic capture, to assess cross-domain generalizability. The experimental results show that unified-trained models reach macro-F1 $\geq 0.948$ across both source domains, decisively outperforming single-source baselines (macro-F1 $\leq 0.74$). Deployment on HospNet26 reveals a significant laboratory-to-clinic covariate shift---$90$--$96\%$ of benign clinical traffic is misclassified as attack, primarily as DoS---providing the first quantitative measurement of this gap and motivating concrete recalibration strategies for safe operational deployment of IoMT intrusion detection systems.

\bigskip
\textbf{Keywords:} Internet of Medical Things, Intrusion Detection System, Machine Learning, Dataset Integration, Network Traffic Analysis, Feature Extraction, Cybersecurity.

\cleardoublepage

% ── Resum (Catalan) ──
\chapter*{Resum}
\addcontentsline{toc}{chapter}{Resum}

La creixent adopció de dispositius de l'Internet de les Coses Mèdiques (\gls{iomt}) en entorns sanitaris ha introduït reptes significatius de ciberseguretat. Aquests dispositius---des de sensors vestibles fins a bombes d'infusió intel·ligents---transmeten dades sensibles de pacients a través de xarxes cada vegada més amenaçades per adversaris. Tot i que els sistemes de detecció d'intrusions (\gls{ids}) basats en aprenentatge automàtic (\gls{ml}) han demostrat potencial, la seva eficàcia està fonamentalment limitada per la qualitat i representativitat de les dades d'entrenament.

Una limitació crítica en la investigació actual és la dependència de conjunts de dades individuals i aïllats que capturen tan sols una porció reduïda del panorama d'amenaces. Models entrenats en aquestes dades tendeixen a sobreajustar-se a topologies de xarxa específiques i fallen en generalitzar a entorns reals. Aquesta tesi aborda aquesta mancança proposant i executant una metodologia sistemàtica per integrar múltiples conjunts de dades heterogenis de tràfic de xarxa \gls{iomt} en un únic conjunt de dades unificat i preparat per a \gls{ml}. La metodologia es basa en l'eina MIOTTA-NPT, que extreu un conjunt estandarditzat de característiques estadístiques directament de captures en brut de paquets, resolent la heterogeneïtat de característiques a l'origen. El procés es complementa amb harmonització d'etiquetes a tres nivells de granularitat, assegurament de la qualitat de les dades, escalat de característiques i seguiment de la procedència.

El conjunt de dades unificat resultant (8.011.534 registres de finestra procedents de CICIoMT2024 i IoMT-TrafficData) s'ha publicat al dipòsit de dades de recerca CORA. Tres classificadors basats en arbres---Decision Tree, Random Forest i XGBoost---van ser entrenats sota règims de font única i unificats, i avaluats en particions de test i en HospNet26, una captura de tràfic d'una xarxa hospitalària real, per quantificar la generalització entre dominis. Els resultats experimentals mostren que els models entrenats amb el conjunt unificat assoleixen macro-F1 $\geq 0{,}948$ en ambdós dominis font, superant decisivament els models de font única (macro-F1 $\leq 0{,}74$). El desplegament en HospNet26 revela un desplaçament de covariable significatiu entre laboratori i clínica---el $90$--$96\%$ del tràfic clínic benigne és classificat erròniament com a atac---proporcionant la primera mesura quantitativa d'aquesta bretxa i motivant estratègies de recalibració per al desplegament operacional segur de sistemes \gls{ids} en entorns hospitalaris.

\bigskip
\textbf{Paraules clau:} Internet de les Coses Mèdiques, Detecció d'intrusions, Aprenentatge automàtic, Integració de dades, Anàlisi de tràfic de xarxa, Ciberseguretat, Generalització entre dominis.

\cleardoublepage

% ── Resumen (Spanish) ──
\chapter*{Resumen}
\addcontentsline{toc}{chapter}{Resumen}

La creciente adopción de dispositivos del Internet de las Cosas Médicas (\gls{iomt}) en entornos sanitarios ha introducido importantes desafíos de ciberseguridad. Estos dispositivos---desde sensores corporales hasta bombas de infusión inteligentes---transmiten datos sensibles de pacientes a través de redes cada vez más amenazadas por adversarios. Aunque los sistemas de detección de intrusiones (\gls{ids}) basados en aprendizaje automático (\gls{ml}) han demostrado potencial, su eficacia está fundamentalmente limitada por la calidad y representatividad de los datos de entrenamiento.

Una limitación crítica en la investigación actual es la dependencia de conjuntos de datos individuales y aislados que capturan únicamente una pequeña parte del panorama de amenazas. Los modelos entrenados en estos datos tienden a sobreajustarse a topologías de red específicas y fallan al generalizarse a entornos reales. Esta tesis aborda esta carencia proponiendo y ejecutando una metodología sistemática para integrar múltiples conjuntos de datos heterogéneos de tráfico de red \gls{iomt} en un único conjunto de datos unificado y preparado para \gls{ml}. La metodología se basa en la herramienta MIOTTA-NPT, que extrae un conjunto estandarizado de características estadísticas directamente desde capturas brutas de paquetes, resolviendo la heterogeneidad de características en el origen. El proceso se complementa con armonización de etiquetas en tres niveles de granularidad, garantía de calidad de los datos, escalado de características y seguimiento de la procedencia.

El conjunto de datos unificado resultante (8.011.534 registros de ventana procedentes de CICIoMT2024 e IoMT-TrafficData) ha sido publicado en el repositorio de datos de investigación CORA. Tres clasificadores basados en árboles---Decision Tree, Random Forest y XGBoost---fueron entrenados bajo regímenes de fuente única y unificados, y evaluados en particiones de test y en HospNet26, una captura de tráfico de una red hospitalaria real, para cuantificar la generalización entre dominios. Los resultados experimentales muestran que los modelos entrenados con el conjunto unificado alcanzan macro-F1 $\geq 0{,}948$ en ambos dominios fuente, superando de forma decisiva a los modelos de fuente única (macro-F1 $\leq 0{,}74$). El despliegue en HospNet26 revela un significativo desplazamiento covariable entre laboratorio y clínica---el $90$--$96\%$ del tráfico clínico benigno es clasificado erróneamente como ataque---proporcionando la primera medición cuantitativa de esta brecha y motivando estrategias de recalibración para el despliegue operacional seguro de sistemas \gls{ids} en entornos hospitalarios.

\bigskip
\textbf{Palabras clave:} Internet de las Cosas Médicas, Detección de intrusiones, Aprendizaje automático, Integración de datos, Análisis de tráfico de red, Ciberseguridad, Generalización entre dominios.
